\subsection{The Field Comment}\label{sec:FieldComment}
Attributes/properties\index{field!attribute}\index{property}, constants\index{field!constant} and \lstinline|enum|\index{enum} values are all summarised under \bsq{field}\index{field} here.

Each and every field\index{field} must have a comment, the \bsq{field comment}, describing it. For private fields\index{field} it can be sufficient to refer to the related getter method\index{method!getter method}\footnote{Of course you have to provide a complete documentation of the property\index{property} there; just a line with a \bdq{\{\nameref{sec:TagReturn} …\}}\index{return@"@return} inline tag will not do the job for the getter method\index{method!getter method} in this case.} instead of writing a lengthy comment into the field comment itself; but do not use the \bdq{\nameref{sec:TagSee}}\index{see@"@see} tag instead of the \bdq{\{\nameref{sec:TagLink} …\}}\index{link@"@link} tag:
\begin{lstlisting}
/**
 *  <p>{@summary For the description of the property 'value' refer to
 *  {@link #getValue()}.
 */
private final Value m_Value;

// AVOID!! Do not use @see
/**
 *  @see #getOtherValue()}.
 */
private final Value m_OtherValue;
\end{lstlisting}

If there is no getter method\index{method!getter method} for the field\index{field}, or if it is not private, a proper description is mandatory. One sentence is enough in most cases, but you are not limited to that.

It is always a good idea to describe the valid values for the field\index{field}, its default value, and whether \lstinline|null| is a possible value for a reference. This is a hard requirement for all non-final non-private fields\index{field} – in particular because there should \textit{never} be any non-final non-private fields\index{field} at all!

For a serialisable class\index{class}\index{serialization}, the fields\index{field} that will be serialised should be tagged with \bdq{\nameref{sec:TagSerial}}\index{serial@"@serial} and the appropriate description.

% Elaborate on comments for constants!!!!





Constants (\lstinline|public static final| fields) that are initialised with a literal have to use the \nameref{sec:TagValue} tag in there description. Also \lstinline|private static final| or \lstinline|protected static final| fields that are initialised with a literal should use the \nameref{sec:TagValue}.

This looks like this:
\begin{lstlisting}
/**
 *  The vested system property for the file encoding used by the JVM:
 *  {@value}.
 */
public static final String PROPERTY_FILE_ENCODING = "file.encoding";
\end{lstlisting}

The comments for enum values are nothing else than field comments for constants – in fact, an enum value is exactly that: a \lstinline|public static final| field initialised with an instance of the enum type.


%==============================================================================
% End of file!!
%==============================================================================

